\documentclass[aspectratio=169]{beamer}

\usetheme[progressbar=frametitle]{metropolis}
\metroset{block=fill}
\usecolortheme{default}
\setbeamercolor{background canvas}{bg=white}
\setbeamercolor{normal text}{fg=black}

\title{Gruppo I - Microplastiche}
\subtitle{Analisi Spettrale per la Distinzione della Grana}
\author{Andrea Bellu, Calvin You, Andrii Myzyn}
\date{\today}

\begin{document}

\begin{frame}
    \titlepage
\end{frame}

\begin{frame}{Introduzione}
    Le microplastiche rappresentano una sfida ambientale crescente a causa della loro persistenza e varietà morfologica. 
    \vfill
    \begin{itemize}
        \item \textbf{Problematica:} La caratterizzazione visiva tradizionale è spesso insufficiente per distinguere microplastiche con proprietà chimico-fisiche simili ma granulometria differente.
        \item \textbf{Tecnologia:} L'imaging iperspettrale (\textit{Hyperspectral Imaging - HSI}) permette di acquisire un intero spettro per ogni pixel dell'immagine (ipercubo), combinando risoluzione spaziale e spettrale.
        \item \textbf{Strumentazione:} Utilizzo della camera \textbf{NIREOS HERA}, basata sulla trasformata di Fourier, operante nel range 400-1000 nm.
    \end{itemize}
\end{frame}

\begin{frame}{Obiettivo del Progetto}
    Il fulcro della ricerca è determinare se la firma spettrale sia influenzata dalla dimensione della grana delle microplastiche.
    \vfill
    \begin{block}{Goal Specifici}
        \begin{enumerate}
            \item \textbf{Differenziazione:} Distinguere campioni a \textit{grana grossa} e \textit{grana fina}.
            \item \textbf{Ottimizzazione Sorgente:} Confrontare l'efficacia dell'illuminazione a \textbf{luce bianca} (riflettanza standard) rispetto alla \textbf{luce UV} (fluorescenza/riflettanza specifica).
            \item \textbf{Validazione:} Analizzare la variabilità spettrale tramite calcolo dello spettro medio e deviazione standard.
        \end{enumerate}
    \end{block}
\end{frame}

\begin{frame}{Setup Sperimentale: Hardware}
    \begin{columns}
        \begin{column}{0.5\textwidth}
            \textbf{Camera Iperspettrale HERA:}
            \begin{itemize}
                \item Tecnologia: Interferometro a ritaglio variabile (Fourier-Transform).
                \item Risoluzione Spaziale: 1280x1024 pixel.
                \item Range: VIS-NIR (400-1000 nm).
            \end{itemize}
            \vspace{0.5cm}
            \textbf{Illuminazione:}
            \begin{itemize}
                \item Lampada LED Bianca (Reflectance mode).
                \item Lampada LED UV (Excitation mode).
            \end{itemize}
        \end{column}
        \begin{column}{0.5\textwidth}
            \begin{figure}
                \centering
                \includegraphics[width=1\textwidth]{images/camera_hera.png} 
            \end{figure}
        \end{column}
    \end{columns}
\end{frame}

\begin{frame}{Configurazione delle Misure}
    Per garantire la confrontabilità dei dati, è prevista una fase cruciale di \textbf{normalizzazione} spettrale per eliminare l'influenza della sorgente e dello sfondo.
    
    \vfill
    \begin{itemize}
        \item \textbf{Misura di Riferimento (Background):} Acquisizione del piano bianco (Teflon) per definire la riflettanza unitaria ($100\%$).
        \item \textbf{Parametri di Acquisizione:}
            \begin{itemize}
                \item \textit{Scan Mode M}: Step di 2.5 nm per un bilanciamento tra tempo e risoluzione.
                \item \textit{Exposure}: Ottimizzata per evitare la saturazione (evidenziata in rosso dal software).
            \end{itemize}
    \end{itemize}

    \begin{block}{Piano di Misura (4 Acquisizioni Target)}
        \centering
        \begin{tabular}{lcc}
            \hline
            \textbf{Campione} & \textbf{Lampada Bianca} & \textbf{Lampada UV} \\
            \hline
            Microplastica Grana Grossa & Misura 1 & Misura 2 \\
            Microplastica Grana Fina & Misura 3 & Misura 4 \\
            \hline
        \end{tabular}
    \end{block}
\end{frame}

\end{document}